\QHeading{画图解决问题。}
\par\noindent


\Problem{① 四（1）班参加绘画比赛的有16人，参加书法比赛的有14人，两项都参加的有6人。参加比赛的一共有多少人？}
\par




\MethodTitle{方法一}
\Answer{{\Formula{\begin{aligned}[t]&\phantom{={}}16+14-6\\&=30-6\\&=24\text{（人）}\end{aligned}}}\par 答：参加比赛的一共有24人。}

\MethodTitle{方法二：先求只属于第二类的人数}
\Answer{第一类人数已经包含两类兼有的人，只需再加上仅属于第二类的人。}
\Answer{\Formula{\begin{aligned}[t]&\phantom{={}}16+(14-6)\\
&=16+8\\
&=24\text{（人）}\end{aligned}}}
\Answer{答：参加比赛的一共有24人。}

\MethodTitle{方法三：分成互不重叠的三类}
\Answer{分别计数“只属于第一类”“两类兼有”“只属于第二类”，每人只数一次。}
\Answer{\Formula{\begin{aligned}[t]&\phantom{={}}(16-6)+6+(14-6)\\
&=10+6+8\\
&=16+8\\
&=24\text{（人）}\end{aligned}}}
\Answer{答：参加比赛的一共有24人。}
\Problem{② 四年级同学参加运动会，参加跑步的有45人，参加跳远的有38人，两项都参加的有12人。四年级参加运动会的一共有多少人？}



\MethodTitle{方法一}
\Answer{{\Formula{\begin{aligned}[t]&\phantom{={}}45+38-12\\&=83-12\\&=71\text{（人）}\end{aligned}}}\par 答：参加运动会的一共有71人。}

\MethodTitle{方法二：先求只属于第二类的人数}
\Answer{第一类人数已经包含两类兼有的人，只需再加上仅属于第二类的人。}
\Answer{\Formula{\begin{aligned}[t]&\phantom{={}}45+(38-12)\\
&=45+26\\
&=71\text{（人）}\end{aligned}}}
\Answer{答：参加运动会的一共有71人。}

\MethodTitle{方法三：分成互不重叠的三类}
\Answer{分别计数“只属于第一类”“两类兼有”“只属于第二类”，每人只数一次。}
\Answer{\Formula{\begin{aligned}[t]&\phantom{={}}(45-12)+12+(38-12)\\
&=33+12+26\\
&=45+26\\
&=71\text{（人）}\end{aligned}}}
\Answer{答：参加运动会的一共有71人。}
\Problem{③ 把两张长15厘米的纸条粘贴在一起，粘贴后总长26厘米。重叠部分长多少厘米？}



\MethodTitle{方法一}
\Answer{{\Formula{\begin{aligned}[t]&\phantom{={}}15\times2-26\\&=30-26\\&=4\text{（厘米）}\end{aligned}}}\par 答：重叠部分长4厘米。}

\MethodTitle{方法二：先求第二张伸出去的部分}
\Answer{总长26厘米减去第一张15厘米，得到第二张露在外面的11厘米；第二张全长减这11厘米就是重叠长。}
\Answer{\Formula{\begin{aligned}[t]&\phantom{={}}15-(26-15)\\
&=15-11\\
&=4\text{（厘米）}\end{aligned}}}
\Answer{答：重叠部分长4厘米。}
\Problem{④ 四年级参加元旦文艺演出，跳舞的有34人，合唱的有50人，两项都参加的有13人。一共有多少人参加演出？}



\MethodTitle{方法一}
\Answer{{\Formula{\begin{aligned}[t]&\phantom{={}}34+50-13\\&=84-13\\&=71\text{（人）}\end{aligned}}}\par 答：参加演出的一共有71人。}

\MethodTitle{方法二：先求只属于第二类的人数}
\Answer{第一类人数已经包含两类兼有的人，只需再加上仅属于第二类的人。}
\Answer{\Formula{\begin{aligned}[t]&\phantom{={}}34+(50-13)\\
&=34+37\\
&=71\text{（人）}\end{aligned}}}
\Answer{答：参加演出的一共有71人。}

\MethodTitle{方法三：分成互不重叠的三类}
\Answer{分别计数“只属于第一类”“两类兼有”“只属于第二类”，每人只数一次。}
\Answer{\Formula{\begin{aligned}[t]&\phantom{={}}(34-13)+13+(50-13)\\
&=21+13+37\\
&=34+37\\
&=71\text{（人）}\end{aligned}}}
\Answer{答：参加演出的一共有71人。}
\Problem{⑤ 四（2）班有56人，参加书法班的有26人，参加绘画班的有20人，两种班都不参加的有15人。两种班都参加的有多少人？}



\MethodTitle{方法一}
\Answer{{\Formula{\begin{aligned}[t]&\phantom{={}}26+20-(56-15)\\&=26+20-41\\&=46-41\\&=5\text{（人）}\end{aligned}}}\par 答：两种班都参加的有5人。}

\MethodTitle{方法二：先求只属于第二类的人数}
\Answer{总人数减去第一类人数，再减去两类都不属于的人数，得到仅属于第二类的人数；再从第二类总人数中扣除。}
\Answer{\Formula{\begin{aligned}[t]&\phantom{={}}20-(56-26-15)\\
&=20-(30-15)\\
&=20-15\\
&=5\text{（人）}\end{aligned}}}
\Answer{答：两种班都参加的有5人。}
